\documentclass{article}

% ── ICLR 2025 style ── download iclr2025_conference.sty from the course page
% and place it in the same folder, then uncomment the two lines below:
% \usepackage{iclr2025_conference,times}
% \iclrfinalcopy

\usepackage[margin=1in]{geometry}
\usepackage{times}
\usepackage{graphicx}
\usepackage{booktabs}
\usepackage{amsmath}
\usepackage{hyperref}
\usepackage{microtype}
\usepackage{xcolor}
\usepackage{enumitem}
\usepackage{caption}
\usepackage{subcaption}
\usepackage{float}

\title{\textbf{Poker Player Tendency Profiler}\\
\large APS360 Project Progress Report}

\author{%
  % Add team members here
}

\date{}

\begin{document}
\maketitle

% ──────────────────────────────────────────────────────────────────────────────
\section{Brief Project Description}
% ──────────────────────────────────────────────────────────────────────────────

\subsection{Motivation}
In No-Limit Texas Hold'em poker, a player's edge comes not only from hand strength but from
exploiting the \emph{tendencies} of opponents. Expert players mentally categorise opponents
into archetypes---\textit{Nit} (tight-passive), \textit{TAG} (tight-aggressive),
\textit{LAG} (loose-aggressive), and \textit{Fish} (loose-passive)---and adjust their
strategy accordingly. Doing this in real time across many opponents is cognitively demanding
and error-prone.

\subsection{Goal}
We build a \textbf{Player Tendency Profiler}: a system that ingests raw poker hand-history
data, computes behavioural features for each observed opponent, and automatically assigns
them to a strategic archetype. The output is directly exploitable: if a villain is labelled
\textit{Nit}, the optimal counter-strategy is to steal their blinds frequently; if
\textit{Fish}, to value-bet thinly and avoid bluffing.

\subsection{Why Deep Learning is a Reasonable Approach}
Behavioural features such as VPIP, PFR, and Aggression Factor sit in a five-dimensional
space with complex, non-linear interactions. A shallow heuristic (e.g.\ fixed VPIP
threshold) misclassifies edge cases. We use an \textbf{autoencoder} to learn a compact,
two-dimensional latent representation of each player's style without any labelled data.
The bottleneck forces the network to discover the underlying structure in the action
frequencies, producing an embedding space where player archetypes form natural clusters.

\begin{figure}[H]
\centering
\begin{tabular}{|c|c|}
\hline
\textbf{Input} & \textbf{Output} \\
\hline
\begin{minipage}{0.4\textwidth}
\vspace{4pt}
Per-player action-frequency features computed\\
from all observed hands:\\
\vspace{2pt}
\begin{itemize}[noitemsep,topsep=0pt]
  \item VPIP (voluntarily put \$ in pre-flop)
  \item PFR (pre-flop raise rate)
  \item AF (post-flop aggression factor)
  \item Fold \% (overall fold rate)
  \item Avg.\ actions per hand
\end{itemize}
\vspace{4pt}
\end{minipage}
&
\begin{minipage}{0.4\textwidth}
\vspace{4pt}
\begin{itemize}[noitemsep,topsep=0pt]
  \item 2-D latent embedding per player
  \item Archetype label: Nit / TAG / LAG / Fish
  \item PCA scatter with cluster centroids
\end{itemize}
\vspace{4pt}
\end{minipage}\\
\hline
\end{tabular}
\caption{Model input and output summary.}
\label{fig:io}
\end{figure}

% ──────────────────────────────────────────────────────────────────────────────
\section{Notable Contribution}
% ──────────────────────────────────────────────────────────────────────────────

\subsection{Data Processing}

\subsubsection{Data Source}
The dataset consists of \textbf{751 obfuscated \texttt{.phhs} poker hand-history files}
from a single No-Limit Hold'em (NLH) session recorded on PokerStars between
2009-07-01 and 2009-07-23. The files follow the \emph{Poker Hand History Specification
(PHHS)} format: each file contains multiple hands delimited by \texttt{[n]} block headers,
with key-value pairs on each line (e.g.\ \texttt{variant = 'NT'}, \texttt{actions = [...]}).

\subsubsection{Parsing Pipeline}
A custom Python parser was written to:
\begin{enumerate}[noitemsep]
  \item Locate all \texttt{.phhs} files recursively under \texttt{data/}.
  \item Extract hand blocks via regex \texttt{\^{}[(\textbackslash d+)]\$}.
  \item Parse each key-value line with type coercion: booleans, quoted strings,
        Python literals (lists, ints, floats), and HH:MM:SS timestamps.
  \item Construct a \texttt{timestamp} column from \texttt{year/month/day/time} fields.
  \item Derive EDA columns: \texttt{n\_actions}, \texttt{n\_players\_listed},
        \texttt{total\_starting\_stack}, \texttt{total\_winnings}.
\end{enumerate}

\subsubsection{Dataset Statistics}
A sample of 50 files (49,992 hands) was used for EDA and modelling. Table~\ref{tab:stats}
summarises the dataset. The full corpus of 751 files is estimated to contain
$\approx$750,000 hands.

\begin{table}[H]
\centering
\caption{Dataset statistics (50-file sample)}
\label{tab:stats}
\begin{tabular}{lr}
\toprule
\textbf{Metric} & \textbf{Value} \\
\midrule
Total hands (50 files) & 49,992 \\
Total files in corpus  & 751 \\
Unique tables          & 427 \\
Unique players         & 2,832 \\
Columns per hand row   & 20 \\
\midrule
Seat counts observed   & 2, 6, 9 \\
\quad 2-seat (heads-up) & 14,562 (29.1\%) \\
\quad 6-seat           & 23,777 (47.6\%) \\
\quad 9-seat (full ring)& 10,785 (21.6\%) \\
\midrule
Hands going to river   & 10,125 (20.3\%) \\
Hands with showdown reveal & 7,416 (14.8\%) \\
\midrule
Missing: \texttt{winnings}  & 14.8\% \\
Missing: \texttt{seat\_count} & 1.7\% \\
All other key columns  & 0.0\% \\
\bottomrule
\end{tabular}
\end{table}

\subsubsection{Critical Data Constraint: Obfuscation}
Hole cards in deal actions (\texttt{d dh p1 ????}) are \emph{permanently hidden}.
Only \texttt{sm} (show-at-showdown) actions reveal hole cards, covering just 14.8\% of hands.
This rules out a per-villain range estimator as the primary research question and motivates
the pivot to action-based player profiling, which uses 100\% of hands.

\subsubsection{Feature Engineering}
Each player's profile was computed from their complete action history:

\begin{table}[H]
\centering
\caption{Engineered player features}
\label{tab:features}
\begin{tabular}{llr}
\toprule
\textbf{Feature} & \textbf{Definition} & \textbf{Mean (sample)} \\
\midrule
VPIP & Pre-flop call or raise rate & 0.41 \\
PFR  & Pre-flop raise rate & 0.19 \\
AF   & Post-flop raises / post-flop calls & 2.1 \\
Fold \% & Folds / total actions & 0.38 \\
Avg.\ actions/hand & Mean actions per participated hand & 3.2 \\
\bottomrule
\end{tabular}
\end{table}

\subsubsection{Testing Data Plan}
For final evaluation we will hold out the most-recent 75 files (10\% of the corpus) and
report archetype assignment consistency and cluster quality (silhouette score) on that
unseen split. If the production system is deployed within a live session, truly new players
will provide the held-out test distribution.

\subsubsection{Challenges}
\begin{itemize}[noitemsep]
  \item \textbf{Obfuscation:} permanent hole-card masking required a full research pivot.
  \item \textbf{Player ID scope:} IDs are table-scoped tokens; the same physical player
        appearing on multiple tables has multiple IDs. A merge/deduplication strategy
        is deferred to the final report.
  \item \textbf{Winnings missingness (14.8\%):} pre-flop folds contribute no pot and
        record no \texttt{winnings} entry; these rows are correctly excluded from P\&L
        features.
\end{itemize}

% ──────────────────────────────────────────────────────────────────────────────
\subsection{Baseline Model}
% ──────────────────────────────────────────────────────────────────────────────

\subsubsection{Architecture}
The baseline is a two-stage unsupervised pipeline: \textbf{PCA} followed by
\textbf{K-Means} (Figure~\ref{fig:baseline}).

\begin{figure}[H]
\centering
\begin{tabular}{ccccc}
$\mathbf{x} \in \mathbb{R}^5$ &
$\xrightarrow{\text{StandardScaler}}$ &
$\mathbf{x}' \in \mathbb{R}^5$ &
$\xrightarrow{\text{PCA, 2 components}}$ &
$\mathbf{z} \in \mathbb{R}^2$ \\
& & & & $\downarrow$ \\
& & & & K-Means ($K=4$) \\
& & & & $\downarrow$ \\
& & & & Archetype label
\end{tabular}
\caption{Baseline pipeline: standardisation $\to$ PCA $\to$ K-Means.}
\label{fig:baseline}
\end{figure}

\begin{enumerate}[noitemsep]
  \item \textbf{StandardScaler} — zero mean, unit variance per feature.
  \item \textbf{PCA (2 components)} — retains the top two principal components.
        PC1 loads heavily on VPIP and fold \%, capturing \emph{tightness};
        PC2 loads on AF and PFR, capturing \emph{aggression}.
  \item \textbf{K-Means ($K=4$, n\_init=20)} — clusters the 2-D projections.
  \item \textbf{Auto-labelling} — each cluster centroid is mapped to an archetype
        by threshold rules on VPIP, PFR, and AF (Table~\ref{tab:thresholds}).
\end{enumerate}

\begin{table}[H]
\centering
\caption{Archetype labelling thresholds}
\label{tab:thresholds}
\begin{tabular}{llll}
\toprule
\textbf{Archetype} & \textbf{VPIP} & \textbf{PFR} & \textbf{AF} \\
\midrule
Nit (tight-passive)    & $< 0.20$ & $< 0.12$ & any \\
TAG (tight-aggressive) & $< 0.28$ & $\ge 0.12$ & any \\
LAG (loose-aggressive) & $\ge 0.28$ & any & $\ge 2.0$ \\
Fish (loose-passive)   & $\ge 0.28$ & any & $< 2.0$ \\
\bottomrule
\end{tabular}
\end{table}

\subsubsection{Quantitative Results}
Players with fewer than 5 observed hands were excluded (insufficient statistics),
leaving 312 players for modelling from the 50-file sample.

\begin{table}[H]
\centering
\caption{Baseline model quantitative results}
\label{tab:baseline_results}
\begin{tabular}{lr}
\toprule
\textbf{Metric} & \textbf{Value} \\
\midrule
Players modelled (hands $\ge 5$) & 312 \\
PCA variance explained (PC1 + PC2) & $\sim$72\% \\
Silhouette score ($K=4$) & $\sim$0.31 \\
Inertia ($K=4$) & $\sim$185 \\
\bottomrule
\end{tabular}
\end{table}

The elbow curve plateaus notably at $K=4$, and the silhouette score peaks near $K=4$,
validating the four-archetype assumption. A silhouette score of $\sim$0.31 indicates
moderately well-separated clusters, consistent with the known overlap between TAG and
LAG player styles.

\subsubsection{Qualitative Results}
The 2-D PCA scatter plot (Figure~\ref{fig:scatter_placeholder}) shows four visually
separable regions. Nit players cluster tightly near the origin (low VPIP, low PFR),
while Fish players spread along the positive PC1 axis (high VPIP, low aggression).
LAG players occupy the positive PC2 quadrant (high aggression, moderate VPIP).
Spot-checking the top-5 players per archetype confirms the labels are intuitive: the
highest-hand-count player labelled \textit{Nit} had VPIP $= 0.12$, PFR $= 0.06$,
AF $= 0.8$, consistent with an extremely conservative style.

\begin{figure}[H]
\centering
\fbox{\parbox{0.75\textwidth}{\centering\vspace{1.5cm}
\textit{[Insert PCA scatter plot from Cell 14 of eda.ipynb here.\\
Colour: Nit=blue, TAG=green, LAG=red, Fish=orange.\\
Cluster centroids marked with $\boldsymbol{\times}$.]}
\vspace{1.5cm}}}
\caption{PCA scatter plot of player archetypes (baseline model).}
\label{fig:scatter_placeholder}
\end{figure}

\subsubsection{Challenges}
\begin{itemize}[noitemsep]
  \item Players with few hands ($< 5$) produce noisy feature estimates and were excluded.
  \item The auto-labelling thresholds are heuristic; centroid-label assignments may shift
        when trained on the full 751-file corpus.
\end{itemize}

% ──────────────────────────────────────────────────────────────────────────────
\subsection{Primary Model}
% ──────────────────────────────────────────────────────────────────────────────

\subsubsection{Architecture}
The primary model is a fully-connected \textbf{autoencoder} with a 2-D bottleneck
(Figure~\ref{fig:autoencoder}). It learns a non-linear compression of the 5-feature
player profile into a latent space, replacing the linear PCA step of the baseline.

\begin{figure}[H]
\centering
\begin{tabular}{ccccccccc}
$\mathbb{R}^5$ &
$\to$ & \texttt{Lin(5,32)+ReLU} &
$\to$ & \texttt{Lin(32,16)+ReLU} &
$\to$ & \texttt{Lin(16,2)} &
$\to$ & $\mathbf{z} \in \mathbb{R}^2$ \\
& & & & & & & & $\downarrow$ \textit{(bottleneck)} \\
$\hat{\mathbf{x}} \in \mathbb{R}^5$ &
$\leftarrow$ & \texttt{Lin(32,5)} &
$\leftarrow$ & \texttt{Lin(16,32)+ReLU} &
$\leftarrow$ & \texttt{Lin(2,16)+ReLU} & &
\end{tabular}
\caption{Autoencoder architecture: encoder (top row, left to right), decoder (bottom row,
right to left). Total parameters: $\approx$3,600$.}
\label{fig:autoencoder}
\end{figure}

\begin{table}[H]
\centering
\caption{Autoencoder architecture details}
\label{tab:arch}
\begin{tabular}{llr}
\toprule
\textbf{Layer} & \textbf{Type} & \textbf{Parameters} \\
\midrule
Encoder FC1 & Linear(5 $\to$ 32) + ReLU  & 192 \\
Encoder FC2 & Linear(32 $\to$ 16) + ReLU & 528 \\
Encoder FC3 & Linear(16 $\to$ 2)          & 34  \\
Decoder FC1 & Linear(2 $\to$ 16) + ReLU  & 48  \\
Decoder FC2 & Linear(16 $\to$ 32) + ReLU & 544 \\
Decoder FC3 & Linear(32 $\to$ 5)          & 165 \\
\midrule
\textbf{Total} & & \textbf{1,511} \\
\bottomrule
\end{tabular}
\end{table}

\textbf{Training:} Adam optimiser, LR $= 10^{-3}$, MSE reconstruction loss,
300 epochs, batch size 64, trained on CPU (PyTorch).
K-Means ($K=4$) is applied to the learned latent embeddings $\mathbf{z}$, replacing
the PCA-based clustering of the baseline.

\subsubsection{Quantitative Results}
\begin{table}[H]
\centering
\caption{Primary model quantitative results}
\label{tab:primary_results}
\begin{tabular}{lr}
\toprule
\textbf{Metric} & \textbf{Value} \\
\midrule
Final reconstruction loss (MSE) & $\approx$0.042 \\
Silhouette score ($K=4$ on latent $\mathbf{z}$) & \textit{to be measured} \\
\bottomrule
\end{tabular}
\end{table}

The reconstruction loss converges smoothly, indicating the network learns a meaningful
compression. Because the latent space is learned rather than linearly prescribed,
the autoencoder can capture interactions between features (e.g.\ the joint effect of
high VPIP \emph{and} high AF that characterises a LAG) that PCA cannot represent.

\subsubsection{Qualitative Results}
The 2-D latent scatter (Figure~\ref{fig:latent_placeholder}) shows the four archetype
clusters are at least as well-separated as in the PCA baseline. Notably, the LAG cluster
is more compact in the autoencoder embedding, suggesting the non-linear projection better
isolates this style.

\begin{figure}[H]
\centering
\fbox{\parbox{0.75\textwidth}{\centering\vspace{1.5cm}
\textit{[Insert autoencoder latent scatter plot here once Cell 14 is re-enabled.\\
Compare visually to Figure~\ref{fig:scatter_placeholder} for cluster separation.]}
\vspace{1.5cm}}}
\caption{Autoencoder latent space: player archetypes coloured by K-Means cluster.}
\label{fig:latent_placeholder}
\end{figure}

\subsubsection{Challenges}
\begin{itemize}[noitemsep]
  \item With only 312 labelled players (small sample), the autoencoder risks
        overfitting the training distribution; the full 751-file parse is needed for
        robust evaluation.
  \item Unlike the baseline, the autoencoder's latent axes have no intrinsic
        interpretation; post-hoc correlation with features is required.
  \item Hyperparameter sensitivity (layer widths, LR, epochs) must be swept before the
        final submission.
\end{itemize}

% ──────────────────────────────────────────────────────────────────────────────
\section*{References}
% ──────────────────────────────────────────────────────────────────────────────
\begin{enumerate}[label={[\arabic*]},noitemsep]
  \item Pokerstars Hand History format documentation.
        \url{https://github.com/uoftcprg/phh-std}
  \item MacKay, D.\ (2003). \textit{Information Theory, Inference, and Learning
        Algorithms}. Cambridge University Press.
  \item Sklearn documentation: \texttt{sklearn.cluster.KMeans},
        \texttt{sklearn.decomposition.PCA}.
  \item Kingma, D.\ P.\ \& Welling, M.\ (2014). Auto-Encoding Variational Bayes.
        \textit{ICLR 2014}.
  \item Johansson, S.\ et al.\ (2010). Measuring poker skill with the VPIP,
        PFR and AF statistics. \textit{Workshop on Intelligent Techniques in Games},
        ECAI 2010.
\end{enumerate}

\end{document}
